\section{引言}
\label{sec:introduction}

在 UTXO 账本中，每笔支付都花费先前交易的输出，因此若接收方等待公开执行，支付链的每一跳都会增加一次确认延迟。假设 Alice 向 Bob 支付 100 CAL，而 Bob 在 Alice 的支付执行之前将其转给 Carol。若 Bob 的支付先执行，账本将消费一个尚不存在的输出，缺失的 100 CAL 就需要一个明确且有资金支持的来源，无论 Alice 的交易稍后到达还是永不到达。

在每一跳都等待公开执行是最直接的替代方案，我们的评估将其用作对照。基于证书的系统把接收方的可用性与公开处理分离开来：FastPay 和 Zef 支持低延迟的认证支付，Sui Lutris 则针对不同类型的状态访问结合了广播路径与共识路径~\cite{fastpay,zef,lutris}。支付通道依靠通道流动性实现快速转账~\cite{blitz,donner}。Next 要问的是：原生账本如何为来源尚未执行的认证支付提供资金，其背后由已注册的组织而非通道作为支撑。

三个困难相互耦合。证书在公开登记之前就已流通，因此仅依据公开缺口进行核算会遗漏已存在的承诺。迟到的来源交易不得既向垫付价值的担保方付款，又向已经花掉该价值的接收方付款。客户端可能持有某个来源交易却永不提交。

我们的核心观察是：对缺失来源负责的一方，已经由随输出流转的证书指明。Next 将缺口记在该直接发行方名下，而不是追溯祖先交易；以由备付金支持的签名额度约束每份已签证书的责任；并通过公开决定关闭未解决的义务，而不改变已存储的历史。

经济闭合之后还存在一个记录问题：原始区块中的已赔付输入仍引用被备付金替代的来源输出。本地例外复核与历史重新导出面向已存储的区块：归档副本重新检查输入，报告备付垫付、授权决定与偿还状态，并导出区块。授权历史修复让导出的正文在输入的原始坐标处直接记录备付扣款引用。决策索引能够从原始正文得到同样的经济答案；额外接口是一个可对照原始 commit 验证的修订正文，其修改授权来自副本已执行的公开前缀。账本审计通过已保存的执行证据复核历史记录~\cite{ia-ccf}，Reparo 通过旧交易哈希返回修复数据~\cite{reparo}；Next 授权替换正文，同时保留已有 commit 与后继区块头所引用的完整身份~\cite{comet-blockid}，而无需再次执行支付。

设计包含三个互补的贡献。

\noindent\textbf{C1：支持一轮续付的分层架构。}证书固定支付输出并标识发行方，公开执行决定账本接受哪笔支付。消费组织成员在签名前，以原子步骤锁定输入并为全部 CAL 输出（含找零）预留额度，一轮并行请求即可形成 TXCer。接收方核验并记录输出后，可以立即请求后继，仅携带直接输入的证书，而复制与公开提交在后台继续。表~\ref{tab:model-contract} 列出各阶段建立的保证。

\noindent\textbf{C2：直接责任驱动的来源解耦。}后继消费来源尚未执行的认证输出时，公开执行在同一原子步骤内记录消费与直接发行方义务。后继区块公布来源证书，使原签署者可以依据保留的批准材料重新提交来源。截止期之后，入块且有效的决定从发行方备付金中赔付缺口，迟到来源再予以偿还。我们证明，签名额度覆盖每份证书（包括尚未公开者）的未结 CAL 责任，将重叠的成员预算~\cite{stingray} 扩展至赔付与偿还。

\noindent\textbf{C3：授权历史修复。}赔付后，原始区块仍引用缺失来源，尽管备付金已提供资金。受限变色龙哈希命令把备付扣款引用写入该输入。每次替换由已提交的赔付决定和精确修订任务授权，保持所有者授权、输出及后继引用、完整 BlockID 不变，且仅写入表示状态。原始命令仍作为重放历史。归档副本因此可复核例外，并按已有身份重新导出修订正文；经济闭合则无需等待修复。

在一台 M4 Max 主机上，采用成员同步提交和 NoSync 钱包存储时，
100 跳链到达最后一笔认证到账的时长中位数为 4.707\,s，
跳间等待公共执行时则为 64.479\,s。
同一当前构建的三轮十四服务 P 实验共完成 21,048 笔付款；
赔付和回款均在暂停的适配服务恢复前完成。
独立实验还评估注入网络延迟和单一故障下的运行行为，
以及可选历史表示的本地读取、存储和维护成本。
